% Generated from validated Article Draft Result. Do not hand-edit; revise the structured draft instead.

2020年に本当に変わったのは、modelを一枚岩のartifactとして見るだけでは説明できなくなったことである。Scaling LawsとGPT-3はscaleとfew-shot / in-context behaviorを結び、APIはaccess surfaceを、RAGはexternal knowledgeを、MeenaとBlenderBotはdialogue-specific system designを、ZeROとGShardはtraining systemsを、diffusionやvisual-token系は別のgeneration paradigmを、human feedbackとtoxicity evaluationはpost-training / evaluation layerを可視化した。これらは同じ技術ではないが、model本体の周囲に独立した設計層が増えたという一点で接続する。


一方で、2020年を現在の生成AIsystemの完成年として描くのは誤りである。instruction followingはまだ後年の形に整理されておらず、human-feedback steeringもsummarizationという限定taskでの研究段階だった。retrieval、dialogue、media generation、API accessも別々のresearch / service surfaceとして存在し、single assistant productへ統合されていない。open-weight ecosystemやdeveloper interfaceも、翌年以降ほど明確な年次軸にはなっていなかった。


したがって2020年を理解するための地図は、『最大modelは何だったか』ではなくlayerの分離である。pretraining scale、in-context interaction、service access、external retrieval、conversation-specific evaluation、distributed training、media representation / generation、preference steering、risk evaluationが別々の操作点になった。翌年には、これらのlayerを既存の巨大pretrained modelへ接続し、用途ごとにadaptする設計がさらに前景化する。


SP-2021-Yで扱うのは、この分離したlayerが再利用可能な基盤の周囲で組み直されていく段階である。parameter-efficient adaptation、instruction / prompt-based steering、open-weight / APIの多様化、developer interface、multimodal representation、retrieval / web grounding、より明示的なevaluation / risk framingが強くなる。2020年から2021年への連続性はあるが、2021年の成果を2020年の出来事として先取りしない。2020年は『分岐が揃った年』、2021年は『それらを接続する設計が前景化した年』と区別する。


この見方に立つと、GPT-3は2020年の中心的anchorでありながら、年全体そのものではない。scaleが能力を押し広げた同じ年に、その規模を支えるsystems、外部知識、利用interface、会話、生成media、人間の好み、望ましくない挙動の評価がそれぞれ独立問題として立ち上がった。生成AIが『一つのmodel』から『複数layerを持つsystem』へ変わるための部品が、2020年に揃い始めたのである。
